\documentclass[11pt, aspectratio=169]{beamer}

\usepackage{pgfpages}

\usepackage{helvet}
\usepackage[default]{lato}
\usepackage[T1]{fontenc}
\usepackage[ansinew]{inputenc}
\usepackage{array}


\usepackage{tikz}
\usepackage{verbatim}
\usetikzlibrary{positioning}
\usetikzlibrary{calc}
\usetikzlibrary{arrows}
\usetikzlibrary{decorations.markings}
\usetikzlibrary{shapes.misc}
\usetikzlibrary{matrix,shapes,arrows,fit,tikzmark}
\usepackage{amsmath}
\usepackage{mathpazo}
\usepackage{hyperref}
\usepackage{lipsum}
\usepackage{multimedia}
\usepackage{graphicx}
\usepackage{multirow}
\usepackage{graphicx}
\usepackage{dcolumn}
\usepackage{bbm}
\newcolumntype{d}[0]{D{.}{.}{5}}

\usepackage{changepage}
\usepackage{appendixnumberbeamer}
\newcommand{\beginbackup}{
   \newcounter{framenumbervorappendix}
   \setcounter{framenumbervorappendix}{\value{framenumber}}
   \setbeamertemplate{footline}
   {
     \leavevmode%
     \hline
     box{%
       \begin{beamercolorbox}[wd=\paperwidth,ht=2.25ex,dp=1ex,right]{footlinecolor}%
%         \insertframenumber  \hspace*{2ex} 
       \end{beamercolorbox}}%
     \vskip0pt%
   }
 }
\newcommand{\backupend}{
   \addtocounter{framenumbervorappendix}{-\value{framenumber}}
   \addtocounter{framenumber}{\value{framenumbervorappendix}} 
}


\usepackage{graphicx}
\usepackage[space]{grffile}
\usepackage{booktabs}

% These are my colors -- there are many like them, but these ones are mine.
\definecolor{blue}{RGB}{0,114,178}
\definecolor{red}{RGB}{213,94,0}
\definecolor{yellow}{RGB}{240,228,66}
\definecolor{green}{RGB}{0,158,115}

\hypersetup{
  colorlinks=false,
  linkbordercolor = {white},
  linkcolor = {blue}
}

%% I use a beige off white for my background
\definecolor{MyBackground}{RGB}{255,253,218}

%% Uncomment this if you want to change the background color to something else
%\setbeamercolor{background canvas}{bg=MyBackground}

%% Change the bg color to adjust your transition slide background color!
\newenvironment{transitionframe}{
  \setbeamercolor{background canvas}{bg=yellow}
  \begin{frame}}{
    \end{frame}
}

\setbeamercolor{frametitle}{fg=blue}
\setbeamercolor{title}{fg=black}
%\setbeamertemplate{footline}[frame number]
\setbeamertemplate{navigation symbols}{} 
\setbeamertemplate{itemize items}{-}
\setbeamercolor{itemize item}{fg=blue}
\setbeamercolor{itemize subitem}{fg=blue}
\setbeamercolor{enumerate item}{fg=blue}
\setbeamercolor{enumerate subitem}{fg=blue}
\setbeamercolor{button}{bg=MyBackground,fg=blue,}


\setbeamertemplate{footline}{%
  \leavevmode%
  \ifnum\value{framenumber}>0
    \hfill%
    \usebeamercolor[fg]{page number in head/foot}%
    \raisebox{1.5ex}{\hspace{1em} \insertframenumber}%
    \hspace{2em} % Adjust spacing as needed
  \fi
}
% If you like road maps, rather than having clutter at the top, have a roadmap show up at the end of each section 
% (and after your introduction)
% Uncomment this is if you want the roadmap!
% \AtBeginSection[]
% {
%    \begin{frame}
%        \frametitle{Roadmap of Talk}
%        \tableofcontents[currentsection]
%    \end{frame}
% }
\setbeamercolor{section in toc}{fg=blue}
\setbeamercolor{subsection in toc}{fg=red}
\setbeamersize{text margin left=1em,text margin right=1em} 

\newenvironment{wideitemize}{\itemize\addtolength{\itemsep}{10pt}}{\enditemize}

\title[]{\textcolor{blue}{Tips + Tricks with Beamer for Economists}}
\author{Hyoungchul Kim}
\institute{Wharton UPenn}
%\institute{\small{\begin{tabular}{c c c}
%Author A &&  Paul Goldsmith-Pinkham  \\
%Somewhere Fancy && FRBNY \\ \\
%
%Author C && Author D   \\
%\multicolumn{3}{c}{Somewhere Fancy} \\
%\end{tabular}}}
\date{\today}


\begin{document}

%%% TIKZ STUFF
\tikzset{   
        every picture/.style={remember picture,baseline},
        every node/.style={anchor=base,align=center,outer sep=1.5pt},
        every path/.style={thick},
        }
\newcommand\marktopleft[1]{%
    \tikz[overlay,remember picture] 
        \node (marker-#1-a) at (-.3em,.3em) {};%
}
\newcommand\markbottomright[2]{%
    \tikz[overlay,remember picture] 
        \node (marker-#1-b) at (0em,0em) {};%
}
\tikzstyle{every picture}+=[remember picture] 
\tikzstyle{mybox} =[draw=black, very thick, rectangle, inner sep=10pt, inner ysep=20pt]
\tikzstyle{fancytitle} =[draw=black,fill=red, text=white]
%%%% END TIKZ STUFF

% Title Slide
\begin{frame}
\maketitle
\addtocounter{framenumber}{-1} % Ensure numbering starts from 1 after title
\end{frame}

\section{Color}
\begin{transitionframe}
  \begin{center}
    { \Huge \textcolor{black}{Color}}
  \end{center}
\end{transitionframe}

\begin{frame}[fragile]{I use two to three colors regularly throughout my presentations}
  \begin{wideitemize}
    \item[-] I use the \textcolor{blue}{blue} color everywhere
    \item[-] I use \textcolor{red}{red}  to contrast and emphasize
    \item[-] I use \textcolor{green}{green}  as a backup
    \item[-] I use \textcolor{yellow}{yellow}  as a transition color
  \end{wideitemize}
\medskip
Make your own color wheel palette!  \url{https://www.sessions.edu/color-calculator/}\\
\end{frame}

\begin{frame}{Background color: I'm undecided}
  \begin{wideitemize}
    \item Changing the background color makes it easier to read
    \item Make your figures in R or Stata have the same background color 
      \begin{itemize}
      \item[] (Or transparent)
      \end{itemize}
    \item In auditoriums \textcolor{red}{large}, go black background with white text \end{wideitemize}
\end{frame}

{  \setbeamercolor{background canvas}{bg=black}
\begin{frame}
  \textcolor{white}{Example for an auditorium -- contrast is much higher}
\end{frame}
}

\begin{frame}[fragile]{Fitting figures doesn't have to be painful}
  \begin{wideitemize}
    \item Simple commands to fit figures: 
    \begin{verbatim}
    \resizebox{0.7\linewidth}{!}{
      \includegraphics{figure1_effect.pdf}}
    \end{verbatim}
    \item The \texttt{\textbackslash linewidth} number is defined within an environment
    \item Simply center with a \texttt{center} environment
  \end{wideitemize}
\end{frame}


\section{Tables}
\begin{transitionframe}
  \begin{center}
    \Huge Tables
  \end{center}
\end{transitionframe}
\begin{frame}[fragile]{Fitting tables doesn't have to be painful}
  \begin{wideitemize}
    \item Simple commands to fit table and center: 
    \begin{verbatim}
\makebox[\linewidth][c]{
\begin{tabular}{l cc ddd}
  results...
\end{tabular}
    \end{verbatim}
    \item Don't use the \texttt{table} environment
    \item If you don't want to center, use change \texttt{[c]} to \texttt{[l]}
  \end{wideitemize}
\end{frame}

\begin{frame}{Highlight cells using tikz (see source code)}
\makebox[\linewidth][c]{
\begin{tabular}{l cc ddd}
  \toprule
  &  Mean at && \multicolumn{3}{c}{Difference-in-Differences Estimates} \\
  \cmidrule{4-6}
  & $ t=-1 $ &&\multicolumn{1}{c}{1 Year} & \multicolumn{1}{c}{2 Years} & \multicolumn{1}{c}{3 Years} \\
  \cmidrule{2-2} \cmidrule{4-6}
  & \multicolumn{1}{c}{(1)}  &&\multicolumn{1}{c}{(2)} & \multicolumn{1}{c}{(3)} & \multicolumn{1}{c}{(4)} \\
  \cmidrule{2-6}
  Outcome 1 & 2.58 && 0.11 &0.08 &\marktopleft{a1} 0.12\\
  & (2.55) && (0.04) & (0.04) & (0.04)\\ 
  Outcome 2 & 60.90 && -0.73 &-1.13 & -1.58\\
  & (17.02) && (0.10) & (0.11) & (0.12)\markbottomright{a1}{red} \\
  Outcome 3 & 18.98 && 0.77 &1.28 & 1.62\\
  & (6.74) && (0.13) & (0.13) & (0.12)\\ 
  \bottomrule
\end{tabular}
}
\uncover<2->{\tikz[overlay,remember picture,inner sep=1pt]
\node[draw=red,rounded corners,fit=(marker-a1-a.north west) (marker-a1-b.south east)] {};}
\end{frame}

\section{Dummy Frames}
\begin{transitionframe}
  \begin{center}
    \Huge Dummy Frames
  \end{center}
\end{transitionframe}

\begin{frame}{Two Column Frame}
\begin{columns}[T] % align columns
\begin{column}{.58\textwidth}
\color{red}\rule{\linewidth}{4pt}
Column 1
\end{column}%
\hfill%
\begin{column}{.38\textwidth}
\color{blue}\rule{\linewidth}{4pt}
Column 2
\end{column}%
\end{columns}
\end{frame}

\begin{frame}{Consider making your math prettier}
  \begin{wideitemize}
    \item Don't overdo it when putting up equations\\(either regressions or theorems)
    \item Try adding color and text to highlight the relevant formula
    \item Consider (Imbens and Angrist 1994): 
      \begin{equation*}
        \alpha_{g}^{IV} = \left. \underbrace{Cov(Y, g(\textcolor{blue}{Z}))}_{\text{\textcolor{red}{Reduced Form}}} \middle/ \underbrace{Cov(D, g(\textcolor{blue}{Z}))}_{\text{\textcolor{green}{First Stage}}} \right.
      \end{equation*}
  \end{wideitemize}
\end{frame}


\begin{frame}{Almost done!}
  \begin{wideitemize}
  \item Use the \texttt{\textbackslash appendix} command to restart the numbering
  \item The frame counter says this is the last slide, but it's not
  \item (Test it, see you on next page)
  \end{wideitemize}
\end{frame}

\appendix

\begin{frame}[label=appendix_start]{Almost done!}
  \begin{wideitemize}
  \item See, now we're in backup slide land
  \item This is made useful by having links throughtout the talk
  \item Here's a button, which is how I make links \hyperlink{appendix_end}{\beamergotobutton{Next slide}}
  \end{wideitemize}
\end{frame}

\begin{frame}[label=appendix_end]{Use it to intimidate audiences!}
  \begin{wideitemize}
    \item[] Now you can make it clear you've done a shitload of work
      \begin{itemize}
      \item[]  without having to show everything! \hyperlink{appendix_start}{\beamergotobutton{Back}}
      \end{itemize}
    \item[] You label a frame with the \texttt{[label=name]} option, and then point a link to it
    \item[] You can make an object a link using the \texttt{\textbackslash hyperlink\{label\}\{object\}} command
  \end{wideitemize}
\end{frame}



\end{document}
